\section*{Preamble}

This work is the continuation of a workshop paper presented at the STM workshop
in 2017, with title \emph{``Authentic Execution of Distributed Event-Driven
Applications with a Small TCB''} and Job Noorman as leading
author~\cite{noorman:authentic-execution}. Two years later, during the academic
year 2019-2020, I had the opportunity to pursue my master's thesis at KU Leuven
and work on this project under the supervision of Fritz Alder, Jan Tobias
M\"{u}hlberg and Frank Piessens. The thesis was titled \emph{``Securing Smart
Environments with Authentic Execution''}~\cite{scopelliti2020thesis}. When my
Ph.D. trajectory began, we decided to further develop this project in
collaboration with Sepideh Pouyanrad. The result of this work is an article
published in 2023 in the TOPS journal and reported in this chapter.

The core idea of this work is to provide a framework to securely instantiate
distributed applications running within \acsp{TEE}, with a strong focus on the
integrity and authenticity of data and computations. In particular, an important
objective is reducing the attack surface as much as possible, allowing each
individual component to be reviewed and audited such that a verifier could not
only attest that the application is running the expected code, but also (and
more importantly) that each output produced by the application could be
\emph{explained} by its source code, along with the provided inputs. Clearly,
such a property is only possible when the software stack, often referred to as
\ac{TCB}, is small enough to be manually reviewed and possibly formally
verified. Isolating small components from a bigger software stack is one of the
main goals of early \acs{TEE} designs like Intel \acs{SGX} and Sancus, and
indeed we leverage them in this work. Recently, however, the trend has started
to take the opposite direction, where huge pieces of code (e.g., entire virtual
machines) are moved inside the \ac{TCB} with a ``lift-and-shift'' approach, as
we will discuss in \cref{ch:cloud}. Yet, solutions such as the one presented in
this chapter may still be relevant for many use cases, especially cyber-physical
systems like smart homes and precision agriculture, both of which have been
explored in our work.

The paper presented at STM 2017 focused on developing the main design of the
framework, which we used to call \emph{``authentic execution''}. Implementation
and evaluation were based on Sancus and supported its secure I/O feature.
However, the number of use cases that we could support was pretty limited, as
Sancus microcontrollers are highly constrained devices that cannot perform
complex operations. Thus, the main focus of my master's thesis was to extend the
implementation to Intel \acs{SGX} in order to support a wider variety of use
cases. Subsequently, the work done during this Ph.D.\ brought about additional
improvements such as a revised design for secure I/O, support for ARM TrustZone
and software updates, and an end-to-end evaluation based on a smart home
application.

Despite being only a research prototype, the authentic execution framework has
become bigger and bigger over the years, and it now consists of thousands of
lines of code spread across thirty repositories, all of which open
source\footnote{\url{https://github.com/AuthenticExecution}}. Some of our
changes have also been integrated into the Sancus
codebase\footnote{\url{https://github.com/sancus-tee}}. In spite of this large
number of repositories and lines of code, significant effort was made to
simplify the development, deployment and testing of workflows by relying on a
high level of automation via scripts, Docker containers, and \ac{CI} pipelines.
As a result, our framework is accessible to all users who are interested in
experimenting with \acsp{TEE} without requiring strong programming skills or a
security background. Apart from the publications mentioned above, our work was
also presented at FOSDEM 2021 and
CCS'21~\cite{authExecFosdem,scopelliti2021poster}. We also used our framework to
evaluate a scenario for secure sensing in supply
chains~\cite{pennekamp2024securing}.